\section{What Is a Field?}

A field is an assignment of physical data to each point of a region of space,
spacetime, or another underlying domain. Instead of describing a system by
finitely many coordinates \(q^1(t),\ldots,q^n(t)\), field theory describes a
system by functions such as
\[
    \phi(t,\vect{x}), \qquad
    \vect{E}(t,\vect{x}), \qquad
    g_{\mu\nu}(x).
\]
The value of the field may be a number, a vector, a tensor, or a more elaborate
geometric object.

The important conceptual change is that the degrees of freedom are distributed.
At every point \(\vect{x}\), a scalar field \(\phi(t,\vect{x})\) has a value.
The collection of all these values forms the state of the field at time \(t\).

\begin{definitionbox}
A \emph{classical field} is a function or geometric assignment whose value is
specified at each point of a continuous domain. In a nonrelativistic setting the
domain is often space and time separately. In a relativistic setting the domain
is spacetime.
\end{definitionbox}

Classical field theory studies the equations that determine how these
assignments vary and interact.
